% titel: Blatt 7, A6: Oberflächenspannung
% skript: Physik I, Kapitel 8.1.3 und 8.2, Seiten 133--134 und 136--139
% stand: Übungsblatt 7, Aufgabe A6

\begin{werkzeugkasten}[Oberflächenenergie und Druckarbeit]
Eine Grenzfläche einer Flüssigkeit besitzt Oberflächenenergie. Bei konstanter
Oberflächenspannung $\sigma$ ist die Änderung $\mathrm dE_\mathrm{ob}
=\sigma\,\mathrm dA$. Die Einheit ist
$[\sigma]=\si{\joule\per\metre\squared}=\si{\newton\per\metre}$.
Bei einer quasistatischen Volumenänderung kann Druckarbeit mit der Änderung der
Oberflächenenergie verglichen werden. Dabei ist die wirksame Druckgrösse die
Differenz zwischen Innen- und Aussendruck, nicht schon der absolute Innendruck.
Die Anzahl der beteiligten Grenzflächen muss zum betrachteten System passen.

\textit{Skript:} Kap. 8.2 und 8.2.1, S. 136--138, insbesondere Gl. (8.6).
\end{werkzeugkasten}

\begin{werkzeugkasten}[Kräfte an einer Grenzfläche]
Oberflächenspannung lässt sich auch als Kraft pro Länge verstehen. Die Kraft
wirkt entlang der Grenzfläche und senkrecht zur betrachteten Kontaktlinie;
für eine vertikale Bilanz zählt nur ihre vertikale Komponente. Diese hängt
von der Orientierung der Grenzfläche beziehungsweise vom Randwinkel ab.
Kräfte längs der gesamten Kontaktlinie müssen addiert werden. Mehrere
Grenzflächen sind nur zu zählen, wenn sie tatsächlich vorhanden sind.

Bei einem Körper an einer Flüssigkeitsoberfläche können neben der Gewichtskraft
auch Druckkräfte der Flüssigkeit und Oberflächenkräfte wirken. Ihr Gleichgewicht
muss mit einem einheitlich definierten Modell formuliert werden. Die Grösse
$m=\rho V$ verbindet dabei Dichte, Volumen und Gewichtskraft $mg$.

\textit{Skript:} Kap. 8.1.3, S. 133--134; Kap. 8.2, S. 137--139,
insbesondere Experiment 8.10 und die Kraft an einer Kontaktlinie.
\end{werkzeugkasten}

\begin{achtung}[Darauf achten]
Eine ideale Oberflächenkraft und ein wirklicher Schwimmzustand sind nicht
identisch: Benetzung, Kontaktwinkel und Verformung der Flüssigkeitsoberfläche
können die Bilanz ändern. Druckdifferenz und absoluter Druck sind verschiedene
Grössen; die Umgebung legt den Bezug fest.
\end{achtung}
